% ============================================================
%  Potenciação — Apostila Premium | PratiqueJá
%  Engine: XeLaTeX
% ============================================================
\documentclass[12pt]{article}
\usepackage{../../pratiqueja}

% \hlm — destaque de resultado em modo-math (cor sem bold)
\newcommand{\hlm}[1]{{\color{darkblue}#1}}

\setsubject{Potenciação}

\begin{document}
\pagenumbering{arabic}
\setstretch{1.2}
\color{bodytext}

\heroheader{Potenciação}%
           {Propriedades, Casos Especiais e Notação Científica}%
           {Matemática · Intermediário}

\setlength{\columnsep}{18pt}
\setlength{\columnseprule}{0.5pt}
\def\columnseprulecolor{\color{exborder}}

\begin{multicols}{2}

%─────────────────────────────────────────────────────────────
\TW{Def}{badgepurple}{Definição}{Multiplicação sucessiva de um número}

\regra{A \textbf{potenciação} é a multiplicação sucessiva de um número
por ele mesmo. $a$ = \textbf{base}, $n$ = \textbf{expoente},
resultado = \textbf{potência}.}

\formula{$a^n = \underbrace{a \cdot a \cdots a}_{n\ \text{vezes}}$}

\exemp{%
  \textbf{Numéricos:}\\
  $4^2 = 16$ {\footnotesize(ao quadrado)} \quad $4^3 = 64$ {\footnotesize(ao cubo)}\\
  $2^5 = 32$\\[3pt]
  \textbf{Algébricos:}\\
  $(x-2)^2 = (x-2)(x-2)$ \quad $(3x)^2 = 9x^2$%
}

%─────────────────────────────────────────────────────────────
\TW{$\pm$}{badgegreen}{Regra dos Sinais na Base}{$-a^n$ é diferente de $(-a)^n$}

\regra{\textbf{Base negativa:} depende da paridade do expoente.\\[2pt]
Expoente \textbf{par} → \textbf{positivo}; \ ímpar → \textbf{negativo}.}

\exemp{%
  $(-2)^4 = +16$ {\footnotesize(par → positivo)}\\[2pt]
  $(-2)^3 = -8$ {\footnotesize(ímpar → negativo)}%
}

\nota{\textbf{O sinal fora do parêntese não é da base:}\\
$-2^4 = -(2^4) = -16$ \quad mas \quad $(-2)^4 = +16$\\
$-5^0 = -1$ \quad mas \quad $(-5)^0 = 1$}

%─────────────────────────────────────────────────────────────
\TW{MB}{darkblue}{Mesma Base}{Multiplicação e divisão}

\SH{Multiplicação — soma os expoentes}
\formula{$a^x \cdot a^y = a^{\,x+y}$}
\exemp{%
  $2^2 \cdot 2^3 = 2^5 = 32$ \qquad $2^5 \cdot 2^{-3} = 2^2 = 4$\\[2pt]
  $x^3 \cdot x^4 = x^7$ \qquad $3x^2 \cdot 2x^3 = 6x^5$%
}

\SH{Divisão — subtrai os expoentes}
\formula{$\dfrac{a^x}{a^y} = a^{\,x-y}$}
\exemp{%
  $\dfrac{2^6}{2^2} = 2^4 = 16$ \qquad
  $\dfrac{2^2}{2^{-3}} = 2^5$ \qquad
  $\dfrac{x^5}{x^2} = x^3$%
}

%─────────────────────────────────────────────────────────────
\TW{Pp}{darkblue}{Produto e Potência de Potência}{Distribuir e compor expoentes}

\SH{Potência de um produto — distribui}
\formula{$(a\cdot b)^n = a^n \cdot b^n$}
\exemp{%
  $(2\cdot3)^4 = 2^4\cdot3^4 = 1296$\\[2pt]
  $(2x)^3 = 8x^3$ \qquad $(3x^2)^2 = 9x^4$%
}

\SH{Potência de potência — multiplica}
\formula{$(a^x)^y = a^{\,x\cdot y}$}
\exemp{%
  $(2^3)^2 = 2^6 = 64$ \qquad $(x^2)^3 = x^6$ \qquad $(5^2)^4 = 5^8$%
}

%─────────────────────────────────────────────────────────────
\TW{FG}{badgepurple}{Fração, Negativo e Zero}{Bases fracionárias e expoentes especiais}

\SH{Potência de fração — distribui}
\formula{$\left(\dfrac{a}{b}\right)^n = \dfrac{a^n}{b^n}$}
\exemp{%
  $\left(\dfrac{5}{3}\right)^2 = \dfrac{25}{9}$ \qquad
  $\left(\dfrac{2}{3}\right)^3 = \dfrac{8}{27}$%
}

\SH{Expoente negativo — inverte a base}
\formula{$a^{-n} = \dfrac{1}{a^n} \qquad \left(\dfrac{a}{b}\right)^{-n} = \left(\dfrac{b}{a}\right)^{n}$}
\exemp{%
  $2^{-3} = \dfrac{1}{8}$ \qquad
  $\left(\dfrac{2}{3}\right)^{-2} = \dfrac{9}{4}$ \qquad
  $10^{-3} = 0{,}001$%
}

\SH{Expoente zero}
\formula{$a^0 = 1 \quad (a \neq 0)$}
\exemp{%
  $7^0 = 1$ \qquad $(-5)^0 = 1$ \qquad $\left(\dfrac{7}{3}\right)^0 = 1$%
}

%─────────────────────────────────────────────────────────────
\TW{!}{vered}{Erro Clássico e Exercício}{Armadilha frequente e regras encadeadas}

\nota{\textbf{A potência NÃO se distribui sobre soma/subtração:}
$(a+b)^n \neq a^n + b^n$ — o erro mais comum.}

\exemp{%
  \textbf{\nok} $(2+3)^2 \neq 2^2 + 3^2 = 13$\\
  \textbf{\ok} $(2+3)^2 = 5^2 = \hlm{25}$\\[3pt]
  \textbf{\nok} $(x+y)^2 \neq x^2 + y^2$\\
  \textbf{\ok} $(x+y)^2 = \hlm{x^2 + 2xy + y^2}$%
}

\SH{Exercício combinado}
\exemp{%
  Simplifique $\dfrac{(2^3)^2 \cdot 2^{-1}}{2^4}$:\\[2pt]
  $(2^3)^2 = 2^6$ {\footnotesize(pot.\ de pot.)}\\
  $2^6 \cdot 2^{-1} = 2^5$ {\footnotesize(mesma base)}\\
  $\dfrac{2^5}{2^4} = 2^1 = \hlm{2}$%
}

%─────────────────────────────────────────────────────────────
\TW{NC}{darkblue}{Notação Científica}{Potências de 10}

\regra{Forma $a \times 10^n$ com $1 \leq |a| < 10$. Expoente positivo →
número grande; negativo → número pequeno.}

\exemp{%
  $6\,000\,000 = 6 \times 10^6$\\[2pt]
  $0{,}00045 = 4{,}5 \times 10^{-4}$\\[2pt]
  $3\,850\,000\,000 = 3{,}85 \times 10^9$%
}

\exemp{%
  $(3\times10^4)(2\times10^3) = 6\times10^7$\\[2pt]
  $\dfrac{8\times10^6}{2\times10^2} = 4\times10^4$%
}

\end{multicols}

%─────────────────────────────────────────────────────────────
\vspace{4pt}
\TW{Ref}{badgegreen}{Principais Potências}{Tabela de referência — quadrados, cubos e quartas}

\vspace{6pt}
\begin{center}
\setlength{\tabcolsep}{14pt}\renewcommand{\arraystretch}{1.3}\small
\begin{tabular}{cccc}
\rowcolor{rulebg}
\textbf{Quadrados} & \textbf{Quadrados} & \textbf{Cubos} & \textbf{Quartas}\\
$1^2 = 1$   & $11^2 = 121$ & $1^3 = 1$    & $1^4 = 1$\\
\rowcolor{exbg}
$2^2 = 4$   & $12^2 = 144$ & $2^3 = 8$    & $2^4 = 16$\\
$3^2 = 9$   & $13^2 = 169$ & $3^3 = 27$   & $3^4 = 81$\\
\rowcolor{exbg}
$4^2 = 16$  & $14^2 = 196$ & $4^3 = 64$   & $4^4 = 256$\\
$5^2 = 25$  & $15^2 = 225$ & $5^3 = 125$  & $5^4 = 625$\\
\rowcolor{exbg}
$6^2 = 36$  & $16^2 = 256$ & $6^3 = 216$  & $6^4 = 1\,296$\\
$7^2 = 49$  & $17^2 = 289$ & $7^3 = 343$  & $7^4 = 2\,401$\\
\rowcolor{exbg}
$8^2 = 64$  & $18^2 = 324$ & $8^3 = 512$  & $8^4 = 4\,096$\\
$9^2 = 81$  & $19^2 = 361$ & $9^3 = 729$  & $9^4 = 6\,561$\\
\rowcolor{exbg}
$10^2 = 100$& $20^2 = 400$ & $10^3 = 1\,000$ & $10^4 = 10\,000$\\
\end{tabular}
\end{center}

\end{document}
